% Created 2025-04-21 lun 17:09
% Intended LaTeX compiler: pdflatex
\documentclass[11pt]{article}
\usepackage[utf8]{inputenc}
\usepackage[T1]{fontenc}
\usepackage{graphicx}
\usepackage{longtable}
\usepackage{wrapfig}
\usepackage{rotating}
\usepackage[normalem]{ulem}
\usepackage{amsmath}
\usepackage{amssymb}
\usepackage{capt-of}
\usepackage{hyperref}
\usepackage{minted}

\usepackage{parskip}
\usepackage[utf8]{inputenc} %% For unicode chars
\usepackage{placeins}
\usepackage[margin=2.50cm]{geometry}
\usepackage[T1]{fontenc}
\usepackage{mathpazo}
\linespread{1.05}
\usepackage[scaled]{helvet}
\usepackage{courier}
\hypersetup{colorlinks=true,linkcolor=blue}
\RequirePackage{fancyvrb}
\DefineVerbatimEnvironment{verbatim}{Verbatim}{fontsize=\small,formatcom = {\color[rgb]{0.5,0,0}}}
\usepackage{mdframed}
\BeforeBeginEnvironment{minted}{\begin{mdframed}}
\AfterEndEnvironment{minted}{\end{mdframed}}
\setcounter{secnumdepth}{3}
\date{}
\title{}
\hypersetup{
 pdfauthor={},
 pdftitle={},
 pdfkeywords={},
 pdfsubject={},
 pdfcreator={Emacs 27.1 (Org mode 9.6.18)}, 
 pdflang={Spanish}}
\begin{document}

\setcounter{tocdepth}{3}
\tableofcontents

\setlength\parindent{10pt}


\section{Introducción a las series temporales y el aprendizaje profundo}
\label{sec:org1295288}
Las series temporales son una secuencia de datos que se han
recolectado a lo largo del tiempo. Estos datos pueden ser medidas de
temperatura, precios de acciones, número de ventas, o cualquier otro
tipo de dato que se recolecte en intervalos de tiempo regulares. Las
series temporales son diferentes a otros tipos de datos porque su
orden y tiempo son esenciales para el análisis. El objetivo del
análisis de series temporales es encontrar patrones y tendencias en
los datos que puedan ser utilizados para hacer predicciones precisas
sobre el futuro.

En los últimos años, el aprendizaje profundo ha revolucionado la forma
en que se analizan las series temporales. Las redes neuronales
recurrentes (RNN) son una arquitectura de aprendizaje profundo que se
utiliza ampliamente en el análisis de series temporales. Las RNN son
capaces de procesar datos secuenciales y utilizar la información
anterior para influir en las predicciones futuras.

Una de las variantes más populares de las RNN son las redes LSTM (Long
Short-Term Memory), que se han utilizado con éxito en el análisis de
series de tiempo debido a su capacidad para aprender a largo plazo y
mantener la información relevante a lo largo del tiempo.

En el siguiente apartado, nos centraremos en las redes LSTM y cómo se
pueden aplicar al análisis de la calidad del aire en series
temporales.

\texttt{=================}

El análisis de series de tiempo es una técnica estadística que se
ocupa de datos que se recopilan a intervalos regulares en el
tiempo. El objetivo del análisis de series de tiempo es extraer ideas
y patrones significativos de los datos, que se pueden utilizar para la
toma de decisiones y la predicción.

El aprendizaje profundo es un subconjunto del aprendizaje automático
que utiliza redes neuronales con múltiples capas para aprender
relaciones complejas entre entradas y salidas. Estas redes neuronales
son capaces de aprender y modelar relaciones no lineales, lo que las
hace adecuadas para el análisis de series de tiempo.

Un tipo popular de red neuronal utilizado para el análisis de series
de tiempo es la red neuronal recurrente (RNN). Las RNN están diseñadas
para manejar datos secuenciales y pueden tener en cuenta las
dependencias temporales entre las observaciones. Las redes LSTM
(Memoria a corto y largo plazo) son un tipo específico de RNN que han
demostrado ser especialmente efectivas en la captura de dependencias a
largo plazo en los datos de series de tiempo.

Al combinar el poder del aprendizaje profundo con las ideas
proporcionadas por el análisis de series de tiempo, es posible
construir modelos potentes que pueden predecir con precisión las
tendencias futuras y tomar decisiones informadas basadas en datos
históricos.


\subsection{LSTM}
\label{sec:org3b77349}
Una de las variantes más eficientes y populares de las redes
neuronales recurrentes (RNN) son las redes LSTM (Long Short-Term
Memory). Las LSTM fueron introducidas para resolver algunos de los
desafíos de las RNN estándar, como la dificultad de aprender
dependencias a largo plazo en datos secuenciales.

Las redes LSTM cuentan con una arquitectura única que incluye
compuertas de entrada, olvido y salida, así como un estado de celda
que permite almacenar información relevante durante periodos
prolongados. Este diseño avanzado les otorga la capacidad de aprender
patrones complejos en datos de series temporales, lo que las hace
ideales para aplicaciones como el procesamiento de lenguaje natural,
la predicción financiera y el análisis de secuencias en general.

A pesar de que las LSTM tienen un número mayor de parámetros en
comparación con las GRU, su capacidad para manejar dependencias más
largas y su flexibilidad han hecho que sean ampliamente adoptadas en
una variedad de tareas donde la precisión es crítica.

En el siguiente apartado, nos centraremos en las redes LSTM y cómo se
pueden aplicar al análisis de la calidad del aire en series
temporales.
\subsubsection{Arquitectura redes LSTM}
\label{sec:orgc1abb98}

Las redes LSTM (Long Short-Term Memory) son un tipo especial de redes
neuronales recurrentes diseñadas para manejar dependencias a largo
plazo en datos secuenciales. Su estructura incluye componentes clave
que trabajan juntos para procesar y almacenar información de manera
eficiente:

\begin{enumerate}
\item \textbf{\textbf{Celdas de memoria}}: Son el núcleo de las LSTM y se encargan de
almacenar información relevante durante largos periodos de
tiempo. Estas celdas permiten que la red decida qué información
conservar o descartar.

\item \textbf{\textbf{Puertas}}:
\begin{itemize}
\item \textbf{\textbf{Puerta de entrada}}: Controla qué información nueva se almacena
en la celda de memoria.
\item \textbf{\textbf{Puerta de olvido}}: Decide qué información existente en la
celda debe ser eliminada.
\item \textbf{\textbf{Puerta de salida}}: Determina qué parte de la información
almacenada en la celda se utiliza para la salida actual.
\end{itemize}

\item \textbf{\textbf{Estados}}:
\begin{itemize}
\item \textbf{\textbf{Estado de celda}}: Es una línea de flujo que transporta
información a lo largo de la red, con interacciones mínimas.
\item \textbf{\textbf{Estado oculto}}: Proporciona información al siguiente paso de
tiempo y a la salida.
\end{itemize}
\end{enumerate}

Estas características hacen que las LSTM sean ideales para tareas como
procesamiento de lenguaje natural, reconocimiento de voz y predicción
de series temporales.

\subsubsection{Funcionamiento red LSTM}
\label{sec:orge671825}
Las redes LSTM (Long Short-Term Memory) funcionan gestionando el flujo
de información a través de pasos temporales en datos secuenciales. A
continuación te explico su funcionamiento de manera más clara:

\begin{enumerate}
\item \textbf{\textbf{Procesamiento inicial}}: En cada paso de tiempo, la LSTM recibe
la entrada actual, el estado oculto y el estado de celda del paso
anterior.

\item \textbf{\textbf{Puerta de olvido}}: Determina si se debe eliminar información del
estado de celda. Esto se hace a través de una función sigmoide que
evalúa la relevancia de cada elemento del estado anterior.

\item \textbf{\textbf{Puerta de entrada}}: Decide qué nueva información se debe agregar
al estado de celda. Aplica una función sigmoide para filtrar la
entrada y una función tangente hiperbólica para generar valores
candidatos.

\item \textbf{\textbf{Actualización del estado de celda}}: El estado de celda se
actualiza combinando la información conservada de la puerta de
olvido y la nueva información de la puerta de entrada.

\item \textbf{\textbf{Puerta de salida}}: Utiliza el estado de celda actualizado para
calcular el estado oculto que se enviará al siguiente paso temporal
y se usará como salida. También utiliza una función sigmoide para
filtrar información relevante.
\end{enumerate}

Gracias a esta arquitectura, las redes LSTM son capaces de recordar
información importante durante largos periodos y descartar datos
irrelevantes. Esto las hace ideales para tareas como traducción
automática, análisis de sentimientos y predicción de series
temporales.

\subsubsection{Preparación de los datos}
\label{sec:org8967c37}
Una preparación adecuada de los datos es esencial para cualquier
análisis de series temporales. Aquí hay algunos pasos que se pueden
tomar para preparar los datos de series temporales para el análisis:

\begin{itemize}
\item Manejo de datos faltantes: Los datos de series temporales a menudo
contienen valores faltantes, que pueden ocurrir debido a varias
razones como fallas en el equipo o errores de entrada de datos. Es
importante manejar los datos faltantes adecuadamente, ya que puede
afectar la precisión del análisis. Los métodos comunes para manejar
datos faltantes incluyen la interpolación, la imputación y la
eliminación.

\item Transformación de datos: Los datos de series temporales pueden
necesitar ser transformados para hacerlos estacionarios o para
mejorar su normalidad. Las transformaciones de datos comunes
incluyen la toma del logaritmo, la raíz cuadrada o la diferencia de
los datos.

\item Normalización: La normalización es el proceso de escalar los datos a
un rango común. Esto se hace a menudo para asegurarse de que todas
las variables tengan igual peso en el análisis. Las técnicas de
normalización comunes incluyen la normalización mín-máx y la
normalización z-score.

\item Selección de características: La selección de características es el
proceso de seleccionar variables relevantes para el análisis. Esto
se puede hacer utilizando conocimiento del dominio o utilizando
técnicas estadísticas como el análisis de correlación o el análisis
de componentes principales.
\end{itemize}

\subsubsection{Creación y entrenamiento de un red LSTM}
\label{sec:org136e07c}
Breve guía sobre cómo crear y entrenar una red LSTM para el análisis
de series temporales:

\begin{enumerate}
\item \textbf{\textbf{Configurar la arquitectura de la red}}: El primer paso para crear
una red LSTM es configurar la arquitectura de la red. Esto implica
especificar el número de capas, el número de neuronas en cada capa
y el tipo de función de activación a utilizar. La arquitectura de
la red dependerá del problema específico que se esté abordando y
puede requerir algo de experimentación para encontrar la
configuración óptima.

\item \textbf{\textbf{Seleccionar hiperparámetros}}: Los hiperparámetros son parámetros
que no son aprendidos por el modelo, sino que son establecidos por
el usuario antes del entrenamiento. Los hiperparámetros comunes
para las redes LSTM incluyen la tasa de aprendizaje, el tamaño del
lote (batch size) y el número de épocas. Estos hiperparámetros
pueden tener un impacto significativo en el rendimiento del modelo
y pueden requerir cierto ajuste para encontrar los valores óptimos.

\item \textbf{\textbf{Preparar los datos}}: Antes de entrenar la red LSTM, los datos
deben prepararse en un formato que pueda alimentarse en la
red. Esto implica dividir los datos en conjuntos de entrenamiento y
prueba, normalizar los datos y remodelarlos en un tensor 3D con
forma (muestras, pasos de tiempo, características).

\item \textbf{\textbf{Entrenar el modelo}}: Una vez que se han especificado la
arquitectura de la red y los hiperparámetros, y se han preparado
los datos, la red LSTM puede entrenarse usando la retropropagación
a través del tiempo (BPTT). Esto implica alimentar los datos de
entrenamiento en la red y actualizar los pesos utilizando el
descenso de gradiente para minimizar la función de pérdida.
\end{enumerate}

Siguiendo estos pasos, se puede crear y entrenar una red LSTM para el
análisis de series temporales. Es importante tener en cuenta que la
construcción de un modelo LSTM efectivo a menudo requiere cierta
experimentación y ajuste fino para lograr resultados óptimos.

\begin{itemize}
\item Ejemplo en Python con Keras
\label{sec:orgbde5daa}
Aquí hay un ejemplo de cómo usar una red LSTM para el análisis de
series de tiempo en Python:
\begin{minted}[]{python}
import numpy as np
from keras.models import Sequential
from keras.layers import LSTM, Dense

# preparar datos
data = np.random.rand(1000, 1) # generar datos aleatorios
timesteps = 10 # número de pasos de tiempo para usar en la entrada
X = np.zeros((data.shape[0]-timesteps, timesteps, data.shape[1]))
y = np.zeros((data.shape[0]-timesteps, data.shape[1]))
for i in range(timesteps, data.shape[0]):
    X[i-timesteps] = data[i-timesteps:i]
    y[i-timesteps] = data[i]

# dividir en conjuntos de entrenamiento y prueba
train_size = int(X.shape[0] * 0.8)
X_train, X_test = X[:train_size], X[train_size:]
y_train, y_test = y[:train_size], y[train_size:]

# crear red LSTM
model = Sequential()
model.add(LSTM(16, input_shape=(timesteps, data.shape[1])))
model.add(Dense(data.shape[1]))
model.compile(loss='mean_squared_error', optimizer='adam')

# entrenar el modelo
model.fit(X_train, y_train, epochs=100, batch_size=32)

# evaluar el modelo
mse = model.evaluate(X_test, y_test)
print('MSE de prueba:', mse)
\end{minted}

\item Ejemplo en Pytorch
\label{sec:org301b06f}
\begin{minted}[]{python}
import torch
import torch.nn as nn
from torch.utils.data import DataLoader, TensorDataset

# Datos ficticios de entrada (batch_size, time_steps, features)
# 100 muestras, 10 pasos de tiempo, 8 características
x = torch.randn(100, 10, 8)
# Salida con 1 característica
y = torch.randn(100, 1)

dataset = TensorDataset(x, y)
dataloader = DataLoader(dataset, batch_size=16, shuffle=True)

class LSTMModel(nn.Module):
   def __init__(self):
       super().__init__()
       self.lstm = nn.LSTM(input_size=8, hidden_size=64, batch_first=True)
       self.fc = nn.Linear(64, 1)

   def forward(self, x):
       out, _ = self.lstm(x)
       # Tomamos el último paso de tiempo
       return self.fc(out[:, -1, :])

model = LSTMModel()
criterion = nn.MSELoss()
optimizer = torch.optim.Adam(model.parameters(), lr=0.001)

# Entrenamiento
for epoch in range(10):
    for batch_x, batch_y in dataloader:
        optimizer.zero_grad()
        output = model(batch_x)
        loss = criterion(output, batch_y)
        loss.backward()
        optimizer.step()
    print(f"Epoch {epoch + 1}, Loss: {loss.item():.4f}")
\end{minted}


Los ejemplos crean una red LSTM simple con una capa LSTM y una capa de
salida densa. La red se entrena con datos generados aleatoriamente y
se evalúa en un conjunto de prueba. El error cuadrático medio (MSE) se
calcula para medir el rendimiento del modelo.

Este es solo un ejemplo simple para ilustrar cómo se puede utilizar
una red LSTM para el análisis de series de tiempo en Python. En la
práctica, la preparación de datos y la arquitectura del modelo
tendrían que adaptarse al problema específico que se está abordando.
\end{itemize}

\subsubsection{Evaluación del modelo}
\label{sec:org89eb3a0}
Vamos a ver, brevemente, cómo evaluar un modelo LSTM para análisis de
series de tiempo:

\begin{enumerate}
\item \textbf{\textbf{Visualización de resultados}}: Una forma de evaluar el
rendimiento de un modelo LSTM es visualizar los valores predichos
frente a los valores reales. Esto se puede hacer usando un gráfico
de líneas, con los valores reales representados en un color y los
valores predichos representados en otro color. Esto le permite
evaluar visualmente qué tan bien el modelo está capturando las
tendencias y patrones en los datos.

\item \textbf{\textbf{Cálculo de métricas de rendimiento}}: Otra forma de evaluar el
rendimiento de un modelo LSTM es calcular métricas de rendimiento
como el error cuadrático medio (MSE) o el error absoluto medio
(MAE). Estas métricas proporcionan una medida cuantitativa de qué
tan bien el modelo está funcionando, con valores más bajos que
indican un mejor rendimiento.

\item \textbf{\textbf{Comparación con modelos de referencia}}: También puede ser útil
comparar el rendimiento del modelo LSTM con modelos de referencia
como un pronóstico ingenuo o un promedio móvil simple. Esto
proporciona un punto de referencia para evaluar el rendimiento del
modelo LSTM y puede ayudar a determinar si la complejidad adicional
del modelo LSTM está justificada.
\end{enumerate}

Siguiendo estos pasos, se puede evaluar el rendimiento de un modelo
LSTM en el análisis de series de tiempo. Es importante tener en cuenta
que diferentes problemas pueden requerir diferentes métodos y métricas
de evaluación, por lo que es importante elegir los métodos adecuados
para el problema específico que se está abordando.

\begin{itemize}
\item Código para evaluar el modelo
\label{sec:org89679bf}
Ejemplo de cómo evaluar un modelo LSTM para el análisis de series
temporales en Python, basado en el ejemplo anterior:
\begin{minted}[]{python}
import matplotlib.pyplot as plt

# generar predicciones
y_pred = model.predict(X_test)

# visualizar resultados
plt.plot(y_test, label='Actual')
plt.plot(y_pred, label='Predicted')
plt.legend()
plt.show()

# calcular métricas de rendimiento
mse = np.mean((y_test - y_pred)**2)
mae = np.mean(np.abs(y_test - y_pred))
print('MSE:', mse)
print('MAE:', mae)

# comparar con el modelo base (pronóstico ingenuo)
y_naive = y_test[:-1]
mse_naive = np.mean((y_test[1:] - y_naive)**2)
print('Naive MSE:', mse_naive)
\end{minted}

\begin{minted}[]{python}
import matplotlib.pyplot as plt
import torch
import numpy as np

# Generar predicciones
model.eval()  # Cambiar el modelo al modo de evaluación
with torch.no_grad():
    # Convertimos las predicciones en formato NumPy
    y_pred = model(X_test).numpy()
    # Aseguramos que los datos reales estén en formato NumPy
    y_test = y_test.numpy()

# Visualizar resultados
plt.plot(y_test, label='Actual')
plt.plot(y_pred, label='Predicted')
plt.legend()
plt.show()

# Calcular métricas de rendimiento
mse = np.mean((y_test - y_pred) ** 2)
mae = np.mean(np.abs(y_test - y_pred))
print('MSE:', mse)
print('MAE:', mae)

# Comparar con el modelo base (pronóstico ingenuo)
y_naive = y_test[:-1]
mse_naive = np.mean((y_test[1:] - y_naive) ** 2)
print('Naive MSE:', mse_naive)
\end{minted}


El código genera predicciones utilizando el modelo LSTM entrenado y
visualiza los resultados utilizando un gráfico de líneas. Se calcula
el error cuadrático medio (MSE) y el error absoluto medio (MAE) para
proporcionar una medida cuantitativa del rendimiento del modelo. El
rendimiento del modelo LSTM también se compara con un pronóstico
ingenuo, que sirve como modelo base.

Es solo un ejemplo simple para ilustrar cómo se puede evaluar un
modelo LSTM para el análisis de series temporales en Python. En la
práctica, los métodos y métricas de evaluación deberían adaptarse al
problema específico que se está abordando.
\end{itemize}


\subsubsection{Mejorar el modelo}
\label{sec:org2f9c122}
Cómo mejorar y ajustar un modelo LSTM para el análisis de series temporales:

\begin{itemize}
\item Optimización de hiperparámetros: Una forma de mejorar el
rendimiento de un modelo LSTM es optimizar los
hiperparámetros. Esto implica probar sistemáticamente diferentes
combinaciones de hiperparámetros para encontrar los valores que
resulten en el mejor rendimiento. Los hiperparámetros comunes para
optimizar incluyen la tasa de aprendizaje, el tamaño del lote y el
número de épocas.

\item Explorar diferentes arquitecturas de red: Otra forma de mejorar el
rendimiento de un modelo LSTM es explorar diferentes arquitecturas
de red. Esto podría implicar agregar o eliminar capas, cambiar el
número de neuronas en cada capa o usar diferentes tipos de capas,
como capas convolucionales o de abandono.

\item Regularización: La regularización es una técnica utilizada para
evitar el sobreajuste mediante la adición de un término de
penalización a la función de pérdida. Las técnicas de
regularización comunes para las redes LSTM incluyen la
regularización L1 y L2, el abandono y la detención temprana.
\end{itemize}

Siguiendo estos pasos, se puede mejorar y ajustar un modelo LSTM para
obtener un mejor rendimiento en el análisis de series temporales. Es
importante tener en cuenta que mejorar un modelo LSTM a menudo
requiere cierta experimentación y ajuste fino para lograr resultados
óptimos.

\begin{itemize}
\item Código
\label{sec:orgea57ded}
ejemplo de cómo  mejorar un modelo LSTM para el  análisis de series de
tiempo en Python, basado en los ejemplos anteriores:
\begin{minted}[]{python}
from keras.layers import Dropout
from keras.callbacks import EarlyStopping

# crear una red LSTM con dropout
model = Sequential()
model.add(LSTM(16, input_shape=(timesteps, data.shape[1])))
model.add(Dropout(0.2))
model.add(Dense(data.shape[1]))
model.compile(loss='mean_squared_error', optimizer='adam')

# entrenar el modelo con early stopping
early_stopping = EarlyStopping(monitor='val_loss', patience=10)
history = model.fit(X_train, y_train, epochs=100, batch_size=32, validation_split=0.2, callbacks=[early_stopping])

# evaluar el modelo
mse = model.evaluate(X_test, y_test)
print('Test MSE:', mse)
\end{minted}

Creamos una red LSTM con una capa de dropout para prevenir el
sobreajuste. La red se entrena usando early stopping para detener el
entrenamiento cuando la pérdida de validación deja de mejorar. El
rendimiento del modelo LSTM mejorado se evalúa en el conjunto de
prueba y se compara con los resultados anteriores.

Este es solo un ejemplo simple para ilustrar cómo se puede mejorar un
modelo LSTM para el análisis de series de tiempo en Python. En la
práctica, mejorar un modelo LSTM a menudo requiere experimentación y
ajustes finos para lograr resultados óptimos.
\end{itemize}

\subsubsection{Siguientes pasos}
\label{sec:orgb90102e}
Aquí hay una breve guía sobre cómo sacar conclusiones e identificar
los próximos pasos después de construir y evaluar un modelo LSTM para
el análisis de series temporales:

\begin{enumerate}
\item \textbf{\textbf{Resumir las conclusiones clave}}: Después de construir y evaluar
un modelo LSTM, es importante resumir las conclusiones clave y las
ideas obtenidas del análisis. Esto podría incluir un resumen del
rendimiento del modelo, cualquier patrón o tendencia interesante
identificada en los datos y cualquier limitación o desafío
encontrado durante el análisis.

\item \textbf{\textbf{Identificar áreas para una exploración adicional}}: La
construcción de un modelo LSTM para el análisis de series
temporales a menudo plantea nuevas preguntas y áreas para una
exploración adicional. Identifique estas áreas y sugiera posibles
próximos pasos para un análisis adicional. Esto podría incluir la
exploración de diferentes fuentes de datos, la prueba de diferentes
arquitecturas de modelo o la aplicación del modelo a nuevos
problemas o dominios.

\item \textbf{\textbf{Comunicar los resultados}}: Finalmente, es importante comunicar
efectivamente los resultados del análisis a las partes
interesadas. Esto podría implicar la creación de visualizaciones o
informes para resumir las conclusiones clave e ideas, y presentar
los resultados de manera clara y concisa.
\end{enumerate}

Siguiendo estos pasos, se pueden sacar conclusiones del análisis e
identificar los próximos pasos para una exploración adicional o la
aplicación del modelo LSTM.
´


dropout 1
Dropout es una técnica de regularización utilizada para prevenir el
sobreajuste en redes neuronales, incluyendo las redes LSTM. Durante el
entrenamiento, el dropout "elimina" o desactiva de manera aleatoria
una proporción de neuronas en una capa, lo que significa que su
contribución al avance y retroceso temporalmente se elimina. Esto
tiene el efecto de reducir la capacidad de la red y alentarla a
aprender representaciones más robustas de los datos.

En una red LSTM, el dropout se puede aplicar a las entradas, salidas o
conexiones recurrentes de la capa LSTM. Aplicar dropout a las entradas
y salidas se puede hacer utilizando la capa Dropout estándar en la
mayoría de los marcos de aprendizaje profundo. Aplicar dropout a las
conexiones recurrentes requiere el uso de una implementación
especializada de la capa LSTM que admite dropout recurrente.

El Dropout es una técnica poderosa para mejorar el rendimiento de las
redes LSTM y prevenir el sobreajuste. Sin embargo, es importante
ajustar cuidadosamente la tasa de dropout para equilibrar el
compromiso entre la capacidad del modelo y la generalización.



\subsection{GRU}
\label{sec:orgaa97e76}


Una de las variantes más eficientes y populares de las redes neuronales recurrentes (RNN) son las redes GRU (Gated Recurrent Unit). Las GRU fueron introducidas como una alternativa más simple a las redes LSTM, pero manteniendo muchas de sus ventajas, como la capacidad para aprender dependencias a largo plazo en datos secuenciales.

Las redes GRU simplifican la arquitectura de las LSTM combinando las compuertas de entrada y olvido en una sola compuerta de actualización, y eliminando el estado de celda separado. Esto reduce la cantidad de parámetros del modelo, lo que puede traducirse en tiempos de entrenamiento más cortos y un rendimiento comparable, e incluso superior, en ciertos conjuntos de datos.

Gracias a su diseño más compacto, las redes GRU han demostrado ser especialmente útiles en situaciones donde se dispone de menos datos o cuando el tiempo de entrenamiento es crítico. También son ampliamente utilizadas en tareas de procesamiento de lenguaje natural, series temporales financieras, y análisis de secuencias en general.



\subsubsection{Estructura y función de las redes GRU}
\label{sec:orgd7cd378}
Las GRU son una variante simplificada de las LSTM que también están diseñadas para manejar dependencias a largo plazo en datos secuenciales. Su estructura incluye:

\begin{enumerate}
\item \textbf{\textbf{Puertas}}:
\begin{itemize}
\item \textbf{\textbf{Puerta de actualización}}: Combina las funciones de las puertas de entrada y olvido de las LSTM. Decide cuánto de la información pasada debe mantenerse y cuánto de la nueva información debe incorporarse.
\item \textbf{\textbf{Puerta de reinicio}}: Controla cuánto de la información pasada se olvida antes de calcular el nuevo estado.
\end{itemize}

\item \textbf{\textbf{Estado oculto}}:
\begin{itemize}
\item En las GRU, no hay un estado de celda separado como en las LSTM. En su lugar, el estado oculto actúa como el único portador de información a lo largo de los pasos temporales.
\end{itemize}
\end{enumerate}

Esta simplificación reduce el número de parámetros en comparación con las LSTM, lo que puede hacer que las GRU sean más rápidas de entrenar y menos propensas al sobreajuste en ciertos casos.

---
\subsubsection{Funcionamiento de las redes GRU}
\label{sec:orge00a051}
El flujo de información en una GRU se gestiona de la siguiente manera:

\begin{enumerate}
\item \textbf{\textbf{Procesamiento inicial}}: En cada paso de tiempo, la GRU recibe la entrada actual y el estado oculto del paso anterior.

\item \textbf{\textbf{Cálculo de la puerta de actualización}}:
\begin{itemize}
\item La puerta de actualización utiliza una función sigmoide para determinar cuánto de la información pasada debe mantenerse y cuánto de la nueva información debe incorporarse.
\end{itemize}

\item \textbf{\textbf{Cálculo de la puerta de reinicio}}:
\begin{itemize}
\item La puerta de reinicio decide cuánto de la información pasada debe olvidarse antes de calcular el nuevo estado. Esto se hace aplicando una función sigmoide.
\end{itemize}

\item \textbf{\textbf{Cálculo del estado candidato}}:
\begin{itemize}
\item Se genera un estado candidato utilizando la entrada actual y el estado oculto modificado por la puerta de reinicio. Esto se realiza mediante una función tangente hiperbólica.
\end{itemize}

\item \textbf{\textbf{Actualización del estado oculto}}:
\begin{itemize}
\item El estado oculto final se calcula como una combinación ponderada del estado anterior y el estado candidato, utilizando los valores de la puerta de actualización.
\end{itemize}
\end{enumerate}

Gracias a esta arquitectura más simple, las GRU son eficaces para tareas como procesamiento de lenguaje natural, reconocimiento de voz y predicción de series temporales, al igual que las LSTM, pero con menor complejidad computacional.

En el siguiente apartado, nos centraremos en las redes GRU y cómo se pueden aplicar al análisis de la calidad del aire en series temporales.

---

\subsubsection{Preparación de los datos}
\label{sec:orgb6eb8d6}
La preparación de los datos para una red GRU es similar a la de una red LSTM. Los pasos clave incluyen:

\begin{itemize}
\item \textbf{\textbf{Manejo de datos faltantes}}: Como en otros modelos de series temporales, es crucial tratar los valores perdidos, utilizando métodos como interpolación, imputación o eliminación de registros, para evitar sesgos o errores en el aprendizaje del modelo.

\item \textbf{\textbf{Transformación de datos}}: Las transformaciones como el logaritmo, la raíz cuadrada o la diferenciación pueden ser necesarias para estabilizar la varianza o hacer los datos más estacionarios, lo cual es beneficioso para el modelado.

\item \textbf{\textbf{Normalización}}: Escalar los datos asegura que las diferentes variables tengan igual influencia en el modelo. Las técnicas comunes son la normalización min-max y la estandarización z-score.

\item \textbf{\textbf{Selección de características}}: Es útil seleccionar variables relevantes para reducir la complejidad del modelo y mejorar el rendimiento. Esto puede hacerse con análisis estadístico o conocimiento del dominio.
\end{itemize}

\subsubsection{Creación y entrenamiento de una red GRU}
\label{sec:org60d95f0}
Guía básica para implementar y entrenar una red GRU:

\begin{enumerate}
\item \textbf{\textbf{Diseñar la arquitectura}}: Se define la estructura de la red, el número de capas GRU, neuronas por capa, funciones de activación y capas adicionales. Las GRU suelen requerir menos parámetros que las LSTM, lo que permite arquitecturas más ligeras.

\item \textbf{\textbf{Elegir hiperparámetros}}: Los parámetros como tasa de aprendizaje, tamaño del batch y número de épocas afectan al rendimiento del modelo. Es común usar técnicas como búsqueda en cuadrícula o validación cruzada para ajustarlos.

\item \textbf{\textbf{Preparar los datos}}: Los datos se deben convertir en tensores tridimensionales (muestras, pasos temporales, características). Además, se deben dividir en conjuntos de entrenamiento y prueba.

\item \textbf{\textbf{Entrenar el modelo}}: Con los datos listos y la arquitectura definida, se entrena la red usando descenso de gradiente y retropropagación a través del tiempo.
\end{enumerate}

\subsubsection{Ejemplo en Python con Keras}
\label{sec:org89f79f3}
\begin{minted}[]{python}
import numpy as np
from keras.models import Sequential
from keras.layers import GRU, Dense

# preparar datos
data = np.random.rand(1000, 1)
timesteps = 10
X = np.zeros((data.shape[0]-timesteps, timesteps, data.shape[1]))
y = np.zeros((data.shape[0]-timesteps, data.shape[1]))
for i in range(timesteps, data.shape[0]):
    X[i-timesteps] = data[i-timesteps:i]
    y[i-timesteps] = data[i]

# dividir en entrenamiento y prueba
train_size = int(X.shape[0] * 0.8)
X_train, X_test = X[:train_size], X[train_size:]
y_train, y_test = y[:train_size], y[train_size:]

# crear modelo GRU
model = Sequential()
model.add(GRU(16, input_shape=(timesteps, data.shape[1])))
model.add(Dense(data.shape[1]))
model.compile(loss='mean_squared_error', optimizer='adam')

# entrenar modelo
model.fit(X_train, y_train, epochs=100, batch_size=32)
\end{minted}

\subsubsection{Ejemplo en PyTorch}
\label{sec:org8961d24}
\begin{minted}[]{python}
import torch
import torch.nn as nn
from torch.utils.data import DataLoader, TensorDataset

x = torch.randn(100, 10, 8)
y = torch.randn(100, 1)

dataset = TensorDataset(x, y)
dataloader = DataLoader(dataset, batch_size=16, shuffle=True)

class GRUModel(nn.Module):
    def __init__(self):
        super().__init__()
        self.gru = nn.GRU(input_size=8, hidden_size=64, batch_first=True)
        self.fc = nn.Linear(64, 1)

    def forward(self, x):
        out, _ = self.gru(x)
        return self.fc(out[:, -1, :])

model = GRUModel()
criterion = nn.MSELoss()
optimizer = torch.optim.Adam(model.parameters(), lr=0.001)

for epoch in range(10):
    for batch_x, batch_y in dataloader:
        optimizer.zero_grad()
        output = model(batch_x)
        loss = criterion(output, batch_y)
        loss.backward()
        optimizer.step()
    print(f"Epoch {epoch + 1}, Loss: {loss.item():.4f}")
\end{minted}

\subsubsection{Evaluación del modelo}
\label{sec:orgff38a85}
Se evalúa el modelo GRU mediante:

\begin{enumerate}
\item \textbf{\textbf{Visualización de resultados}}: Comparar las predicciones del modelo con los valores reales mediante gráficos.

\item \textbf{\textbf{Métricas de rendimiento}}: Cálculo de métricas como MSE y MAE para cuantificar el rendimiento.

\item \textbf{\textbf{Comparación con modelos base}}: Comparar el rendimiento del modelo GRU con modelos sencillos como promedio móvil o pronóstico ingenuo.
\end{enumerate}

\subsubsection{Mejorar el modelo}
\label{sec:orga49b008}
Algunas estrategias para optimizar un modelo GRU incluyen:

\begin{itemize}
\item \textbf{\textbf{Tuning de hiperparámetros}}: Ajustar número de unidades GRU, tasa de aprendizaje, número de épocas, etc.

\item \textbf{\textbf{Regularización}}: Uso de dropout para prevenir el sobreajuste.

\item \textbf{\textbf{Arquitecturas alternativas}}: Experimentar con redes GRU bidireccionales o combinarlas con otras capas como convolucionales o de atención.
\end{itemize}

\subsubsection{Código con mejoras}
\label{sec:org8924f6e}
\begin{minted}[]{python}
from keras.layers import Dropout
from keras.callbacks import EarlyStopping

model = Sequential()
model.add(GRU(16, input_shape=(timesteps, data.shape[1])))
model.add(Dropout(0.2))
model.add(Dense(data.shape[1]))
model.compile(loss='mean_squared_error', optimizer='adam')

early_stopping = EarlyStopping(monitor='val_loss', patience=10)
history = model.fit(X_train, y_train, epochs=100, batch_size=32, validation_split=0.2, callbacks=[early_stopping])
\end{minted}

\subsubsection{Siguientes pasos}
\label{sec:org0139d66}
Después de implementar una red GRU:

\begin{enumerate}
\item \textbf{\textbf{Extraer conclusiones}}: Analizar los resultados y rendimiento del modelo.

\item \textbf{\textbf{Identificar mejoras futuras}}: Evaluar posibles ajustes al modelo, nuevas fuentes de datos o tareas relacionadas.

\item \textbf{\textbf{Comunicar resultados}}: Preparar visualizaciones e informes claros para compartir con otros equipos o partes interesadas.
\end{enumerate}

---
\end{document}
